\documentclass{article}
\usepackage[utf8]{inputenc}
\usepackage[T1]{fontenc}
\usepackage{geometry}
\geometry{
    a4paper,
    left=2.5cm,
    right=2.5cm,
    top=2.5cm,
    bottom=2.5cm
}
\usepackage{graphicx}
\usepackage{float}
\usepackage{hyperref}
\usepackage[english]{babel}
\usepackage{amsmath}
\usepackage{amssymb}
\usepackage{array}
\usepackage{booktabs}
\usepackage{caption}
\usepackage{subcaption}
\usepackage{enumitem}
\usepackage{fancyhdr}
\usepackage{titlesec}
\usepackage{url}
\usepackage{textcomp}
\usepackage{xcolor}
\usepackage{listings}
\usepackage{multirow}
\usepackage{longtable}
\usepackage{tabularx}
\usepackage{times}

% Set up page style
\pagestyle{fancy}
\fancyhf{}
\rfoot{\thepage}
\renewcommand{\headrulewidth}{0pt}

% Title formatting
\titleformat{\chapter}[display]{\normalfont\bfseries}{\chaptertitle\thechapter}{0pt}{\Huge}
\titlespacing*{\chapter}{0pt}{50pt}{40pt}

% Document info
\title{Template Document}
\author{Academic Support Agent}
\date{\today}

\begin{document}
\maketitle

%-------
%---CAPA

\begin{titlepage}
\begin{center}
\begin{figure}
    \centering
    \includegraphics[height=!, width= 2cm]{imagens/uepb.png}
\end{figure}

\textbf{UNIVERSIDADE ESTADUAL DA PARA BA\\
CAMPUS X\\
CENTRO XXXXXXXXX\\
DEPARTAMENTO DE XXXXXXXX\\
CURSO DE XXXXXXXXX}

\vspace{4cm}


\textbf{NOME DO ESTUDANTE}   

\vspace{4cm}
\textbf{T TULO DO TRABALHO}

\vfill

\textbf{CIDADE\\
ANO} 
    
\end{center}
\end{titlepage}

%----------------
%--FOLHA DE ROSTO

\thispagestyle{empty}
\begin{center}
 
    
\textbf{NOME DO ESTUDANTE}    

\vspace{5cm}

\textbf{T TULO DO TRABALHO}   

\vspace{5cm}

\begin{flushright}
\parbox{8cm}{
\singlespacing{\begin{hyphenrules}{nohyphenation}
Trabalho de Conclus o de Curso apresentado ao Departamento de XXXXXXX do Centro de XXXXXXXXXX da Universidade Estadual da Para ba como requisito parcial   obten o do t tulo de Licenciado(a) em XXXXXXX.        % natureza do trabalho
\vspace{0.5cm}    

\textbf{ rea de concentra o}:   xxxxxxxxxxxxx
\end{hyphenrules}}}
\end{flushright}


\vspace{2cm}

\textbf{Orientador}: Prof. xxxxxxxxxxx
\vfill


\textbf{CIDADE\\
ANO}
\end{center}
\newpage

%----------------
%--FOLHA DE ROSTO

\thispagestyle{empty}
\begin{center}
\textbf{NOME DO ALUNO}

\vspace{1cm}

\textbf{T TULO DO TRABALHO}

\end{center}

\vspace{.4cm}
\begin{flushright}
\parbox{8cm}{
\singlespacing{\begin{hyphenrules}{nohyphenation}
Trabalho de Conclus o de Curso apresentado ao Departamento de XXXXXXXX do Centro de XXXXXXXX da Universidade Estadual da Para ba como requisito parcial   obten o do t tulo de Licenciado(a) em XXXXXXX.
\vspace{0.5cm}

\textbf{ rea de concentra o}:  xxxxxxxxxx
\end{hyphenrules}}}
\end{flushright}

\vspace{.6cm}

\noindent Aprovado em: ---------

\vspace{1.5cm}

\begin{center}
\textbf{BANCA EXAMINADORA}
\vspace{1.5cm}

\noindent\rule{12cm}{0.4pt} \\
Prof. Dr. xxxxxxxxxxxxxx (Orientador)\\
Universidade Estadual da Para ba (UEPB)
\vspace{2cm}

\noindent\rule{12cm}{0.4pt}\\
Profa. Dra. xxxxxxxxxxxxxx  \\
Universidade Estadual da Para ba (UEPB)
\vspace{2cm}

\noindent\rule{12cm}{0.4pt}\\
Prof. Dr. xxxxxxxxxxxxxxx \\
Universidade XXXXXX (CASO MEMBRO EXTERNO)
\vspace{2cm}

\end{center}

%--------------
%---DEDICAT RIA

\thispagestyle{empty}
\null
\vfill
\begin{flushright}
\parbox{4cm}{
\singlespacing{\begin{hyphenrules}{nohyphenation}
Dedicat ria.
\end{hyphenrules}
}}
\end{flushright}
\thispagestyle{empty}
\newpage

%-----------------
%---AGRADECIMENTOS

\thispagestyle{empty}
\section*{\normalsize \centering \textbf{AGRADECIMENTOS}}

Escreva aqui os seus agradecimentos...

\newpage

%-----------
%---EP GRAFE

\thispagestyle{empty}
\null
\vfill
\begin{flushright}
\parbox{8cm}{
\singlespacing{\begin{hyphenrules}{nohyphenation}
``Ep grafe - Cita o relacionada com o tema do trabalho, com indica o de autoria. Dever seguir as diretrizes da NBR 10.520/2002 de Cita o em documentos.''
\end{hyphenrules}
}}
\end{flushright}
\newpage

%---------
%---RESUMO

\thispagestyle{empty}
\section*{\normalsize \centering \textbf{RESUMO}}

\noindent  O resumo dever  ser uma apresenta o concisa dos pontos relevantes de um documento. Deve ressaltar o objetivo, o m todo, os resultados e as conclus es do documento. Dever  ser composto de uma sequ ncia de frases concisas, afirmativas e n o de enumera o de t picos. Recomenda-se o uso de par grafo  nico. Deve-se usar o verbo na voz ativa e na terceira pessoa do singular.

\vspace*{0.5cm}

\noindent\textbf{Palavras-chave:} Palavras-chave 1. Palavras-chave 2. Palavras-chave 3.
 % no m ximo 4 palavras-chave separadas entre si por ponto
 \newpage

%---------
%---ABSTRACT

\thispagestyle{empty}
\section*{\normalsize \centering \textbf{ABSTRACT}}

\noindent Resumo em ingl s

\vspace*{0.5cm}

\noindent\textbf{Keywords:} Keyword 1. Keyword 2. Keyword 3.
\newpage

%----------
%---SUM RIO



%----------
%---INTRODU O

\chapter{INTRODU O}
\thispagestyle{myheadings}

Texto, texto texto...

Texto, texto texto...

%----------
%---CAP TULO I

\chapter{TEORIA DOS N MEROS}
\thispagestyle{myheadings}

Texto, texto texto...

Texto, texto texto...

\section{Se o prim ria}

Texto, texto texto...

Texto, texto texto...

\subsection{Se o secund ria}

Texto, texto texto...

Texto, texto texto...

%----------
%---CAP TULO II

\chapter{T TULO DO CAP TULO}
\thispagestyle{myheadings}

Texto, texto texto...

Texto, texto texto...

\section{Isto   uma se o}

Texto, texto texto...

Texto, texto texto...

\subsection{Isto   uma subse o}

Texto, texto texto...

Texto, texto texto...

%----------
%---CAP TULO III

\chapter{T TULO DO CAP TULO}
\thispagestyle{myheadings}

Texto, texto texto...

Texto, texto texto...

\section{Isto   uma se o}

Texto, texto texto...

Texto, texto texto...

\subsection{Isto   uma subse o}

Texto, texto texto...

Texto, texto texto...

%--------------
%---CAP TULO VI

\chapter{CONCLUS O}
\thispagestyle{myheadings}

Texto, texto texto...

Texto, texto texto...

\section{Isto   uma se o}

Texto, texto texto...

Texto, texto texto...

\subsection{Isto   uma subse o}

Texto, texto texto...

Texto, texto texto...

%-----------------------------
%---REFER NCIAS BIBLIOGR FICAS

\addcontentsline{toc}{chapter}{\protect\numberline{}REFER NCIAS}
\begin{thebibliography}{20}
\thispagestyle{myheadings}

\bibitem[Alencar Filho(1981)]{alencarfilho1981}
ALENCAR FILHO, Edgard de. \textbf{Teoria elementar dos n meros}. S o Paulo: Nobel, 1981.
exemplo de cita o: (ALENCAR FILHO, 1981)
exemplo de cita o: Alencar Filho, 1981

Inserir exemplo de refer ncia de Artigo





% exemplo de refer ncia

\end{thebibliography}

%------------
%---AP NDICES

\addcontentsline{toc}{chapter}{\protect\numberline{}AP NDICE A -- T TULO DO AP NDICE}
\addtocontents{toc}{\protect\setcounter{tocdepth}{0}}

\setcounter{chapter}{0}
\renewcommand{\thechapter}{\Alph{chapter}}
\chapter*{\hfill AP NDICE A -- T TULO DO AP NDICE \hfill}
\stepcounter{chapter}
\thispagestyle{myheadings}

Conte do do Ap ndice A...

Conte do do Ap ndice A...
\end{document}